\documentclass[11pt,a4paper]{article}
\usepackage[margin=1in]{geometry}
\usepackage{parskip}
\usepackage{hyperref}
\usepackage{tabularx}
\usepackage{booktabs}
\begin{document}
\begin{center}
{\LARGE \textbf{AI-assisted Portfolio Project Report}}\\[0.5em]
{\large Ahmed Elkhodary (900213472)}\\
Course: Design of Web-based Systems \\ 
Supervisor: Dr Islam Tharwat \\
\end{center}

\section*{1. Project Overview}
\textbf{Objective:} The objective of this project was to create a professional online CV and portfolio with interactive features and GitHub API integration, using AI coding tools to assist in development.

\textbf{Features Implemented:}
\begin{itemize}
  \item HTML CV structure and semantic sections (About, Experience, Projects, Skills, Contact).
  \item Responsive CSS styling with light/dark themes and accessible components.
  \item JavaScript interactivity: navigation highlighting, section toggles, project detail expansion, dynamic skills addition, contact form validation, and a print-based "Download CV" action.
  \item GitHub API integration to fetch and display public repositories dynamically.
  \item Documentation and deployment instructions (README, AI log).
\end{itemize}

\section*{2. AI Tool(s) Used}
\begin{tabularx}{\textwidth}{@{}lX@{}}
\toprule
Tool Name & Purpose / Task Assisted \\
\midrule
GitHub Copilot (VS Code) & Generated HTML snippets, CSS patterns, and JavaScript helper functions used throughout the site. \\
OpenAI Copilot Chat / ChatGPT & Reviewed code, suggested refactors, and produced the LaTeX report and additional documentation content. \\
\bottomrule
\end{tabularx}

\section*{3. Tasks Assisted by AI}
\begin{enumerate}
  \item \textbf{HTML/CSS scaffolding}: AI provided initial layout and style suggestions; I edited markup for semantics and accessibility.
  \item \textbf{JavaScript functions}: AI suggested patterns for event handling and API integration; I improved validations and error handling.
  \item \textbf{LaTeX report generation}: AI produced the report template and filled content which I adapted and verified.
\end{enumerate}

\section*{4. Challenges Encountered}
AI-generated suggestions were occasionally incomplete (edge-case handling) or stylistically inconsistent with existing code. Debugging and testing were necessary to ensure the features worked reliably across browsers.

\section*{5. Reflections on AI Usage}
\begin{enumerate}
  \item What worked well: Rapid scaffolding, repetitive code generation, and idea exploration.
  \item What needed manual correction: Accessibility attributes, error handling, and API resilience.
  \item Productivity: AI accelerated initial development, allowing focus on integration and polish.
  \item Understanding: Using AI required careful review which reinforced my knowledge of event-driven JS and responsive CSS.
  \item Future use: I would continue using AI tools as assistants but maintain manual verification and testing.
\end{enumerate}

\section*{6. Ethical and Practical Considerations}
Relying on AI raises concerns about correctness, licensing of generated snippets, and potential over-dependence. Developers must verify generated code, attribute third-party content properly, and avoid using AI as a substitute for critical thinking.

\section*{7. Online Portfolio Link}
Hosted URL: https://____________________________

\section*{8. Conclusion}
This AI-assisted project provided practical experience using code-completion and chat-based AI tools to implement a professional portfolio. The tools improved speed but required hands-on validation for correctness and accessibility. Future improvements include stronger API error handling, adding tests, and deploying the site publicly.

\end{document}
